\section{Introduction}
\label{sec:intro}

Vision-language models (VLMs) have demonstrated remarkable capabilities in understanding and reasoning over visual content, from natural images~\cite{radford2021clip,openai2023gpt4v} to documents~\cite{tang2023unifying} and user interfaces~\cite{you2024ferretui,hong2024cogagent}. A particularly consequential application frontier is the deployment of VLMs as autonomous UI agents that navigate web pages, extract structured information from dashboards, and execute multi-step workflows on behalf of users~\cite{zheng2024gpt4vision}. In these settings, the model must not only produce a correct answer but also ground its evidence in a specific visual region of the interface, since downstream consumers (auditors, compliance teams, or chained agents) need to verify \emph{where} the answer came from.

Real-world user interfaces, however, are not static. A single web application may render differently depending on the user's theme preference (light versus dark mode), browser window size, scroll position, or the presence of collapsible navigation panels. We refer to this natural visual variability as \emph{UI drift}: semantically equivalent content presented under perturbed visual conditions. UI drift is distinct from adversarial perturbation in that it arises from routine usage rather than deliberate attack, yet it can substantially alter the pixel-level appearance of an interface without changing its underlying information content. For a VLM agent to be reliable, it must produce consistent answers and accurate evidence localization regardless of these surface-level variations.

Despite the rapid proliferation of UI understanding benchmarks~\cite{liu2024visualwebbench,hsiao2022screenqa,chen2021websrc,baechler2024screenai}, existing evaluations predominantly test models on canonically rendered screenshots. They measure whether a model can answer a question about a static screenshot but do not assess whether the model would give the same answer, or locate the same evidence region, if that screenshot were rendered under a different theme, viewport size, or layout configuration. This gap leaves a critical blind spot in our understanding of VLM robustness for deployed UI agents.

We draw inspiration from the visual robustness benchmarking tradition in image classification~\cite{hendrycks2019imagenetc,hendrycks2019cifarc}, where controlled corruption taxonomies (noise, blur, weather, digital transformations) have proven invaluable for diagnosing model fragility. Extending this paradigm to the UI understanding setting requires a fundamentally different corruption taxonomy: instead of Gaussian noise and fog, the relevant perturbations are theme switches, layout reflows, and viewport changes that preserve semantic content while altering visual appearance.

In this work, we introduce \bench{}, a controlled micro-benchmark designed to measure VLM robustness to UI drift under grounded retrieval conditions. Our benchmark consists of 50 base UI pages spanning five common UI categories (dashboards, settings panels, data tables, articles, and analytics views), each rendered under five calibrated drift types at increasing severity. The resulting 300 images are paired with 390 grounded question-answer (QA) pairs, where each QA pair includes both a ground-truth answer and a ground-truth evidence bounding box identifying the visual region that supports the answer.

We evaluate three frontier VLMs, namely GPT-4o, GPT-4o-mini, and Gemini~2.0~Flash, and report five metrics: Answer Exact Match (\aem{}), Answer Containment (\ac{}), Evidence Grounding IoU, Answer Consistency, and a composite Drift Robustness Score (\drs{}). Our experiments yield several findings with practical implications for the deployment of VLM-based UI agents:

\begin{itemize}
    \item \textbf{Drift robustness and base accuracy are decoupled.} Gemini~2.0~Flash achieves the highest \drs{} (0.950) despite lower base exact-match accuracy (0.928) than GPT-4o (0.949), demonstrating that models can be simultaneously less accurate on clean inputs and more robust to perturbation.
    
    \item \textbf{Evidence grounding is universally weak.} The highest grounding IoU achieved by any model under any condition is 0.115, and the best model-level average is 0.066 (GPT-4o). This means that even when VLMs produce correct answers, they cannot reliably point to the supporting evidence, a critical limitation for audit-ready and explainable systems.
    
    \item \textbf{Task-specific vulnerability profiles differ across models.} Gemini~2.0~Flash fails on trend detection tasks (0.300 base EM), while GPT-4o-mini struggles with table lookups (0.714 base EM). These complementary failure modes suggest that model selection should be informed by the downstream task distribution.
    
    \item \textbf{Composite drift degrades weaker models disproportionately.} GPT-4o-mini's answer consistency drops to 0.800 under composite drift, while GPT-4o maintains 0.970, indicating that model scale provides meaningful robustness benefits under compounded perturbations.
\end{itemize}

To our knowledge, \bench{} is the first benchmark to systematically evaluate VLM robustness to realistic UI drift under grounded retrieval conditions. We release the full benchmark (images, QA pairs, bounding boxes, and evaluation code) to support future research on robust, grounded UI understanding.
